\pdfbookmark[0]{Notices. Herschel: Cotes’s theorem. 1814.}{anzeigen-herschel-1814}
\begin{center}
\rule{.85\textwidth}{.3pt}\par
Göttingische gelehrte Anzeigen.\quad 14 February 1814.\par
\rule{.85\textwidth}{.3pt}
\end{center}
\hspace*{\parindent}\textit{Philosophical Transactions of the Royal Society of London for the Year} 1813. IV and 304 pages. 26 pages of the \textit{Meteorological Journal} and 8 pages of the \textit{Index}, in quarto.

\textit{Mathematical and astronomical papers. --- On a remarkable Application of} \textsc{Cotes}\textit{'s Theorem, by} \textsc{J. F. W. Herschel} (son of the astronomer). Let \(n\) be any integer, let \(n\omega=360^0\), and let \(N\) be any angle. Under these assumptions, \textsc{Cotes}'s theorem gives a simple expression for the product of all the radius vectors that, in astronomical terms, correspond on a conic section to the true anomalies \(N,\ N+\omega,\ N+2\omega,\ N+3\omega\mathpunct{.}\mathpunct{.}\mathpunct{.}\mathpunct{.}N+(n-1)\omega\). Though these and the other similar developments that make up the paper present no particular difficulties in themselves, the paper is a pleasure to read because of the way the subject is treated. Some more recent mathematical writers use the notation \(\cos^2 A\) for the square of \(\cos A\), entirely contrary to analogy, according to which it ought to mean the cosine of an arc \(=\cos A\). What the author says about this notation has our full approval.

